% Présentation 1 : Profil du candidat
% Carlos Denner dos Santos
% Journée d'évaluation - 23 janvier 2026
% Poste de professeur en gestion de l'intelligence artificielle
% Département SIMQG, École de gestion, Université de Sherbrooke

% === PDF SETTINGS (full font embedding) ===
\pdfminorversion=7
\pdfobjcompresslevel=0

\documentclass[aspectratio=169,11pt]{beamer}

% === ENCODAGE ET LANGUE ===
\usepackage[utf8]{inputenc}
\usepackage[T1]{fontenc}
\usepackage[french]{babel}
\usepackage{csquotes}
\usepackage{lmodern}  % Latin Modern fonts (fully embeddable)

% === STYLE UdeS ===
\usepackage{beamer_usherbrooke}

% === PACKAGES ADDITIONNELS ===
\usepackage{graphicx}
\usepackage{booktabs}
\usepackage{array}
\usepackage{colortbl}
\usepackage{tikz}
\usetikzlibrary{shapes,arrows,positioning,calc}
\usepackage{hyperref}

% === MÉTADONNÉES ===
\title[Profil du candidat]{Profil du candidat}
\subtitle{Poste de professeur en gestion de l'intelligence artificielle}
\author[C. D. dos Santos]{Carlos Denner dos Santos Jr.}
\institute[UdeS - SIMQG]{
  Département des systèmes d'information\\
  et méthodes quantitatives de gestion (SIMQG)\\
  École de gestion, Université de Sherbrooke
}
\date{23 janvier 2026}

\begin{document}

% ============================================
% SLIDE TITRE
% ============================================
\begin{frame}
  % Bannière verte avec titre de la présentation
  \begin{beamercolorbox}[wd=\paperwidth,ht=2.5cm,dp=0.5cm,center]{palette primary}
    {\Large\bfseries Présentation 1 : Profil du candidat}
  \end{beamercolorbox}
  
  \vspace{0.8cm}
  
  \begin{center}
    {\Large Carlos Denner dos Santos Jr.}
    
    \vspace{0.8em}
    
    {\normalsize Département des systèmes d'information\\
    et méthodes quantitatives de gestion (SIMQG)\\
    École de gestion, Université de Sherbrooke}
    
    \vspace{0.8em}
    
    {\normalsize 23 janvier 2026}
    
    \vspace{1em}
    
    {\large\textcolor{ushervert}{\textbf{\guillemotleft~Je transforme l'IA en capacité de gestion.~\guillemotright}}}
  \end{center}
\end{frame}

% ============================================
% PLAN DE LA PRÉSENTATION
% ============================================
\begin{frame}{Plan de la présentation}
  \tableofcontents
\end{frame}

% ============================================
% SECTION 1 : FORMATION ET PARCOURS (~6 min)
% ============================================
\section{Formation et parcours professionnel}

% --- SLIDE 1.1 : Profil en bref ---
\begin{frame}{Profil en bref}
  \begin{columns}[T]
    \begin{column}{0.55\textwidth}
      \begin{block}{Qui suis-je?}
        \begin{itemize}
          \item \emphvert{Chercheur en systèmes d'information}
          \item \emphvert{Professeur universitaire} (15+ ans)
          \item \emphvert{Consultant en science des données et IA}
        \end{itemize}
      \end{block}
      \vspace{0.5em}
      \begin{alertblock}{Fil conducteur}
        \textbf{20 ans} à l'intersection de l'analytique, des systèmes d'information et de la stratégie
      \end{alertblock}
    \end{column}
    \begin{column}{0.42\textwidth}
      \centering
      \begin{block}{Domaines d'expertise}
        \begin{itemize}
          \item \textbf{IA/ML}
          \item \textbf{SI/Données}
          \item \textbf{Gestion}
        \end{itemize}
      \end{block}
    \end{column}
  \end{columns}
\end{frame}

% --- SLIDE 1.2 : Parcours académique (Timeline) ---
\begin{frame}{Parcours académique}
  \small
  \begin{center}
    \begin{tabular}{c|l}
      \textcolor{ushervert}{\textbf{1999}} & Baccalauréat en Administration \\
      \textcolor{ushervert}{\textbf{2005}} & Maîtrise (UFMG -- Brésil) \\
      \textcolor{ushervert}{\textbf{2009}} & \textbf{Ph.D. Systèmes d'information} (SIUC -- USA) \\
      \textcolor{ushervert}{\textbf{2011}} & Postdoc Informatique (USP -- Brésil) \\
      \textcolor{ushervert}{\textbf{2011}} & Postdoc Statistique (Nottingham -- UK) \\
    \end{tabular}
  \end{center}
  
  \vspace{0.5em}
  \begin{block}{Formation multidisciplinaire}
    \centering
    \textbf{Gestion} + \textbf{Informatique} + \textbf{Statistique} = \emphvert{Perspective unique en IA}
  \end{block}
\end{frame}

% --- SLIDE 1.3 : Expérience professionnelle ---
\begin{frame}{Expérience professionnelle}
  \begin{columns}[T]
    \begin{column}{0.48\textwidth}
      \begin{block}{Académique}
        \begin{itemize}
          \item \textbf{Professeur associé}\\Université de Brasilia (2012--présent)\\{\scriptsize\textit{Congé sans solde depuis 2024}}
          \item \textbf{Co-chercheur principal}\\Jooay App, McGill (2020--présent)
        \end{itemize}
      \end{block}
    \end{column}
    \begin{column}{0.48\textwidth}
      \begin{block}{Industrie}
        \begin{itemize}
          \item \textbf{Expert IA / Consultant}\\Videns AI (2025--présent)
          \item \textbf{Consultant science des données}\\Bell Canada (2022--2024)
          \item Secteurs : Télécom, Consulting, Énergie, Secteur public
        \end{itemize}
      \end{block}
    \end{column}
  \end{columns}
  
  \vspace{0.5em}
  \begin{alertblock}{Valeur ajoutée}
    \centering
    Rare combinaison : \emphvert{rigueur académique} + \emphvert{expérience terrain}
  \end{alertblock}
\end{frame}

% ============================================
% SECTION 2 : EXPERTISES ALIGNÉES AU POSTE (~6 min)
% ============================================
\section{Expertises alignées au poste}

% --- Transition vers section 2 ---
\begin{frame}{}
  % Bannière verte de transition
  \begin{beamercolorbox}[wd=\paperwidth,ht=1.8cm,dp=0.3cm,center]{palette primary}
    {\large\bfseries Section 2 : Expertises alignées au poste}
  \end{beamercolorbox}
  
  \vspace{1.5cm}
  
  \centering
  {\Large\textit{Ce parcours multidisciplinaire m'a préparé idéalement\\pour ce poste...}}
  
  \vspace{1em}
  
  {\large\emphvert{Voyons en quoi mon profil répond précisément aux exigences}}
  \vfill
\end{frame}

% --- SLIDE 2.1 : Tableau d'alignement ---
\begin{frame}{Alignement avec les exigences du poste}
  \small
  \begin{tabular}{>{\raggedright}p{4.2cm}|c|>{\raggedright\arraybackslash}p{6.5cm}}
    \rowcolor{ushervert}
    \textcolor{white}{\textbf{Exigences}} & & \textcolor{white}{\textbf{Mon profil}} \\
    \hline
    \rowcolor{ushervertclair}
    \textbf{Gouvernance} \& risques IA & $\checkmark$ & \textbf{541} citations -- article majeur (2021) \\
    \hline
    \textbf{AIOps / MLOps}\textsuperscript{*} & $\checkmark$ & \textbf{20 ans} exp. terrain + projets LLM \\
    \hline
    \rowcolor{ushervertclair}
    \textbf{Impact stratégique} de l'IA & $\checkmark$ & Stratégie IA (Novacap, Laurentide) \\
    \hline
    \textbf{Performance} organisationnelle & $\checkmark$ & Dashboards opérationnels (Bell, \textbf{3 ans}) \\
    \hline
    \rowcolor{ushervertclair}
    \textbf{Sécurité} \& résilience IA & $\checkmark$ & Sécurité LLM\textsuperscript{**}, prompt injection \\
    \hline
    \textbf{Enseignement} TI/IA & $\checkmark$ & \textbf{15+ ans} tous cycles \\
  \end{tabular}
  
  \vspace{0.2em}
  {\tiny \textsuperscript{*}AIOps/MLOps : ingénierie et opérations d'IA (équivalent DevOps pour l'IA)}\\
  {\tiny \textsuperscript{**}LLM : Large Language Models (grands modèles de langage, ex: GPT)}
  
  \vspace{0.3em}
  \begin{center}
    \emphvert{Forte ad\'equation avec les exigences cl\'es -- preuves acad\'emiques et terrain}
    
    \vspace{0.2em}
    
    {\small\textit{Et je suis aussi ici pour apprendre des forces d\'ej\`a pr\'esentes au SIMQG}}
  \end{center}
\end{frame}

% --- SLIDE 2.2 : Publications phares ---
\begin{frame}{Publications phares en gouvernance de l'IA}
  \begin{columns}[T]
    \begin{column}{0.48\textwidth}
      \begin{block}{Ethics \& Info. Tech. (2021)}
        \centering
        \textbf{``AI Regulation: a framework''}
        
        \vspace{0.3em}
        {\Large\textcolor{usheror}{\textbf{541 citations}}}
        
        \vspace{0.3em}
        {\small Cadre intégrateur $\rightarrow$ \emphvert{outil de décision publique}}
      \end{block}
    \end{column}
    \begin{column}{0.48\textwidth}
      \begin{block}{Gov. Info. Quarterly (2025)}
        \centering
        \textbf{``AI governance in practice''}
        
        \vspace{0.3em}
        {\Large\textcolor{ushervert}{\textbf{28 org. / 5 continents}}}
        
        \vspace{0.3em}
        {\small Écart principes$\leftrightarrow$pratiques $\rightarrow$ \emphvert{implications de gouvernance}}
      \end{block}
    \end{column}
  \end{columns}
  
  \vspace{0.5em}
  \begin{alertblock}{Pipeline de publications}
    \textbf{CACM} (R\&R) $\bullet$ \textbf{ICML 2026} $\bullet$ \textbf{AOM 2026} (x2) $\bullet$ \textbf{Economi\textit{A}}
    
    {\tiny CACM : LLM Firewall | ICML : Causal Prompt Interventions | AOM : stratégie \& théorie | Economi\textit{A} : TI \& salaires}
  \end{alertblock}
\end{frame}

% --- SLIDE 2.3 : Recherche actuelle - Sécurité LLM ---
\begin{frame}{Recherche actuelle : Sécurité des systèmes LLM}
  \begin{columns}[T]
    \begin{column}{0.48\textwidth}
      \begin{block}{Thèmes émergents}
        \begin{itemize}
          \item \emphvert{Défense contre la prompt injection}
          \item Détection et atténuation des \emphvert{hallucinations}
          \item Sécurité des \emphvert{agents autonomes}
          \item \emphvert{Boucles affectives} humain-LLM
        \end{itemize}
      \end{block}
      
      \begin{block}{Analyse de brevets (2010--2025)}
        \begin{itemize}
          \item \textbf{250 000+} familles de brevets IA/ML
          \item Identification des domaines émergents
          \item White-space scoring $\rightarrow$ priorités R\&D
        \end{itemize}
      \end{block}
    \end{column}
    \begin{column}{0.48\textwidth}
      \begin{alertblock}{Risques LLM}
        \centering
        \begin{tabular}{c|c}
          \textcolor{red!70}{\textbf{Attaques}} & \textcolor{red!70}{\textbf{Fiabilité}} \\
          \hline
          \textcolor{red!70}{\textbf{Autonomie}} & \textcolor{red!70}{\textbf{Humain}} \\
        \end{tabular}
      \end{alertblock}
      
      \vspace{0.3em}
      
      \begin{exampleblock}{Transfert : LLM Firewall}
        {\small Pare-feu LLM issu de ma recherche}
        \begin{itemize}\small
          \item Prototype fonctionnel
          \item \emphvert{Recherche $\rightarrow$ artefact d\'eployable}
        \end{itemize}
      \end{exampleblock}
    \end{column}
  \end{columns}
\end{frame}

% --- SLIDE 2.4 : Expérience terrain ---
\begin{frame}{Exemples marquants de projets IA menés}
  \small
  \begin{columns}[T]
    \begin{column}{0.32\textwidth}
      \begin{block}{Consulting}
        \textbf{Novacap}
        
        {\footnotesize\textit{Stratégie IA}}
        \begin{itemize}\small
          \item Due diligence IA
          \item Évaluation portefeuille
        \end{itemize}
        \vspace{0.2em}
        \textbf{\href{http://laurentide-dashboard-alb-2099058886.ca-central-1.elb.amazonaws.com/}{Contrôles Laurentide}}
        
        {\footnotesize\textit{SCALE AI -- ML}}
        \begin{itemize}\small
          \item Optimisation inventaire
          \item Déploiement prod.
        \end{itemize}
      \end{block}
    \end{column}
    \begin{column}{0.32\textwidth}
      \begin{block}{Télécom}
        \textbf{Bell Canada}
        
        {\footnotesize\textit{3 ans -- Data Science}}
        \begin{itemize}\small
          \item Migration 5G
          \item Big Data (Hadoop)
          \item Tableaux de bord BI
          \item Optimisation énergie
        \end{itemize}
      \end{block}
    \end{column}
    \begin{column}{0.32\textwidth}
      \begin{block}{Secteur public}
        \textbf{Brésil}
        
        {\footnotesize\textit{Min. Planejamento}}
        \begin{itemize}\small
          \item Software Público Brasileiro
          \item Gouvernance TI
          \item Données ouvertes
        \end{itemize}
      \end{block}
    \end{column}
  \end{columns}
  
  \vspace{0.3em}
  \begin{alertblock}{}
    \centering
    \textbf{Du concept à la production} -- Maîtrise complète du cycle AIOps/MLOps
  \end{alertblock}
\end{frame}

% ============================================
% SECTION 3 : CONTRIBUTIONS ET SYNERGIES (~6 min)
% ============================================
\section{Contributions et synergies institutionnelles}

% --- Transition vers section 3 ---
\begin{frame}{}
  % Bannière verte de transition
  \begin{beamercolorbox}[wd=\paperwidth,ht=1.8cm,dp=0.3cm,center]{palette primary}
    {\large\bfseries Section 3 : Contributions et synergies}
  \end{beamercolorbox}
  
  \vspace{1.5cm}
  
  \centering
  {\Large\textit{Au-delà de l'adéquation de mon profil...}}
  
  \vspace{1em}
  
  {\large\emphvert{Comment puis-je contribuer activement\\au SIMQG et à l'UdeS?}}
  \vfill
\end{frame}

% --- SLIDE 3.1 : Indicateurs de recherche ---
\begin{frame}{Indicateurs de recherche}
  \begin{columns}[T]
    \begin{column}{0.48\textwidth}
      \begin{block}{Production scientifique}
        \begin{center}
          {\Huge\textcolor{ushervert}{\textbf{59}}} {\small articles de revue}
          \hspace{0.3cm}
          {\Huge\textcolor{usheror}{\textbf{1 776}}} {\small citations}
        \end{center}
        \vspace{0.3em}
        \begin{itemize}
          \item Indice h : \textbf{17} | i10-index : \textbf{31}
          \item \textbf{1 172} citations depuis 2021
          \item Top \textbf{10\%} ResearchGate
        \end{itemize}
      \end{block}
    \end{column}
    \begin{column}{0.48\textwidth}
      \begin{block}{Encadrement}
        \begin{itemize}
          \item $\sim$\textbf{9} doctorats dirigés
          \item $\sim$\textbf{6} maîtrises dirigées
          \item \textbf{3} postdoctorats supervisés
          \item \textbf{21} projets de recherche
        \end{itemize}
        \vspace{0.5em}
        \emphvert{Capacité démontrée} de former la relève
      \end{block}
    \end{column}
  \end{columns}
\end{frame}

% --- SLIDE 3.2 : Intégration SIMQG / CROI ---
\begin{frame}{Intégration au SIMQG et synergies}
  \begin{columns}[T]
    \begin{column}{0.48\textwidth}
      \begin{block}{SIMQG}
        \begin{itemize}
          \item Recherche en gouvernance IA
          \item Enseignement GTA / SIA
          \item Encadrement cycles supérieurs
        \end{itemize}
      \end{block}
      \vspace{0.3em}
      \begin{block}{Centre L. Beaudoin}
        \begin{itemize}
          \item Partenariats industrie
          \item Réseau Bell, Cofomo, Videns AI
        \end{itemize}
      \end{block}
    \end{column}
    \begin{column}{0.48\textwidth}
      \begin{block}{CROI}
        \begin{itemize}
          \item Organisations intelligentes
          \item Performance des systèmes IA
        \end{itemize}
      \end{block}
      \vspace{0.3em}
      \begin{block}{École de gestion}
        \begin{itemize}
          \item Microprogramme IA (3\textsuperscript{e} cycle)
          \item École d'été en IA
        \end{itemize}
      \end{block}
    \end{column}
  \end{columns}
  
  \vspace{0.5em}
  \begin{alertblock}{}
    \centering
    \textit{``Je m'investirai pleinement en \textbf{recherche}, \textbf{enseignement}, \textbf{encadrement} et \textbf{service}.''}
  \end{alertblock}
\end{frame}

% --- SLIDE 3.3 : Vision et mission SIMQG ---
\begin{frame}{Alignement avec la mission du SIMQG}
  \begin{block}{Mission du département}
    \centering
    {\large\textit{``Développer, par la recherche et l'enseignement, des connaissances en management de l'information et de l'intelligence.''}}\vspace{0.3em}\\
    {\small -- Mission du SIMQG}
  \end{block}
  
  \vspace{1em}
  \begin{alertblock}{Ma contribution}
    \begin{itemize}
      \item Recherche appliquée en \emphvert{gouvernance de l'IA}
      \item Pont entre \emphvert{théorie} (publications) et \emphvert{pratique} (industrie)
    \end{itemize}
  \end{alertblock}
\end{frame}

% --- SLIDE 3.4 : Financement et collaborations ---
\begin{frame}{Potentiel de financement et collaborations}
  \begin{columns}[T]
    \begin{column}{0.48\textwidth}
      \begin{block}{Sources de financement visées}
        \begin{itemize}
          \item \textbf{CRSH} -- Gouvernance de l'IA
          \item \textbf{FRQSC} -- Équipes et projets
          \item \textbf{Mitacs} -- Partenariats industriels
          \item \textbf{CRSNG} -- Aspects techniques
          \item Contrats de recherche (secteur privé/public)
        \end{itemize}
      \end{block}
    \end{column}
    \begin{column}{0.48\textwidth}
      \begin{block}{Collaborations établies}
        \begin{itemize}
          \item McGill University (Jooay)
          \item Bell Canada
          \item Cofomo / Videns AI
          \item Réseau international établi
        \end{itemize}
        \vspace{0.3em}
        \emphvert{Prêt à mobiliser} ces réseaux au profit de l'UdeS
      \end{block}
    \end{column}
  \end{columns}
\end{frame}

% ============================================
% CONCLUSION
% ============================================
\section*{Conclusion}

% --- Transition vers conclusion ---
\begin{frame}{}
  % Bannière verte de transition
  \begin{beamercolorbox}[wd=\paperwidth,ht=1.8cm,dp=0.3cm,center]{palette primary}
    {\large\bfseries Conclusion : Pourquoi moi?}
  \end{beamercolorbox}
  
  \vspace{1.5cm}
  
  \centering
  {\Large\textit{En résumé, si je devais vous convaincre\\en cinq points...}}
  
  \vspace{1em}
  
  {\large\emphvert{Voici ce qui fait de moi le bon candidat}}
  \vfill
\end{frame}

\begin{frame}{En résumé : Pourquoi moi?}
  \begin{columns}[T]
    \begin{column}{0.6\textwidth}
      \begin{enumerate}
        \item[1.] \emphvert{Formation multidisciplinaire} \\
        {\small Ph.D. SI + postdocs informatique \& statistique}
        \vspace{0.25em}
        \item[2.] \emphvert{Expertise alignée au poste} \\
        {\small Gouvernance IA, AIOps/MLOps, sécurité LLM}
        \vspace{0.25em}
        \item[3.] \emphvert{Publications de haut niveau} \\
        {\small Revues Q1, pipeline actif, impact international}
        \vspace{0.25em}
        \item[4.] \emphvert{Expérience terrain} \\
        {\small Télécom, consulting, secteur public}
        \vspace{0.25em}
        \item[5.] \emphvert{Engagement total} \\
        {\small Recherche, enseignement, \textbf{encadrement}, service}
      \end{enumerate}
    \end{column}
    \begin{column}{0.35\textwidth}
      \centering
      \colorbox{ushervert}{\parbox{3cm}{\centering\textcolor{white}{\large\bfseries Le bon profil\\pour le SIMQG}}}
      
      \vspace{0.6em}
      
      {\small\textbf{Je suis ici pour bâtir la discipline de gestion de l'IA au SIMQG.}}
      
      \vspace{0.5em}
      
      $\checkmark$ \small Prêt dès le jour 1
    \end{column}
  \end{columns}
\end{frame}

% ============================================
% ANNEXE : PREUVES VÉRIFIABLES (backup)
% ============================================
\begin{frame}{Annexe : Références et méthodologie}
  \scriptsize
  \begin{columns}[T]
    \begin{column}{0.48\textwidth}
      \begin{block}{Publications phares}
        \begin{enumerate}
          \item Santos Jr., Denner (2021). ``A framework for regulating AI''. \textit{Ethics and Information Technology}. \textbf{541 citations} (Google Scholar, jan. 2026)
          \item Santos Jr., Denner (2019). ``Analysis of Patent-based Indicators''. \textit{Bras. Bus. Review}. 59 articles/1776 citations analysés.
          \item Santos Jr. et Tumelero (2024). ``AI Readiness in IS''. \textit{En révision}.
        \end{enumerate}
      \end{block}
    \end{column}
    \begin{column}{0.48\textwidth}
      \begin{block}{Méthodologie brevets}
        \begin{itemize}
          \item \textbf{Source} : USPTO PatentsView API
          \item \textbf{Périmètre} : IPC G06N (IA/ML)
          \item \textbf{Déduplication} : Familles DOCDB
          \item \textbf{Embeddings} : OpenAI text-embedding-3
          \item \textbf{Volume} : 250K familles dédupliquées
        \end{itemize}
      \end{block}
      
      \begin{block}{Benchmark latence LLM Firewall}
        \begin{tabular}{lr}
          \textbf{Mesure} & \textbf{Résultat} \\
          \hline
          Inférence locale & 8.2ms \\
          + Réseau (AWS) & 45ms \\
          Baseline (Rebuff) & 120ms \\
        \end{tabular}
      \end{block}
    \end{column}
  \end{columns}
\end{frame}

% --- SLIDE FINAL ---
\begin{frame}{}
  % Bannière verte finale
  \begin{beamercolorbox}[wd=\paperwidth,ht=1.8cm,dp=0.3cm,center]{palette primary}
    {\large\bfseries Merci de votre attention}
  \end{beamercolorbox}
  
  \vspace{1cm}
  
  \centering
  {\Huge\emphvert{Merci!}}
  
  \vspace{0.8em}
  
  {\Large Questions?}
  
  \vspace{1em}
  
  \begin{tabular}{cl}
    Courriel : & carlosdenner@gmail.com \\
    Téléphone : & +1 (438) 836-4116 \\
    LinkedIn : & linkedin.com/in/carlosdenner \\
    GitHub : & github.com/carlosdenner \\
  \end{tabular}
\end{frame}

\end{document}
